\chapter{图表、公式格式}
图、表和公式均按章编号，使用交叉引用命令维护编号。以下给出插图、分图、三线表、续表和公式的写法。
\section{图的格式}
图应清晰、规范且随文出现。图题位于图下，宋体五号；编号与题名之间空两个汉字。图的前后各留一行。引用外部图片时，应在图题右上角标注文献来源。

\begin{figure}[htbp]
  \centering
  \fbox{\parbox[c][3cm][c]{0.7\linewidth}{\centering 此处替换为论文插图}}
  \caption{插图位置示例}\label{fig:placeholder}
\end{figure}

如图\ref{fig:placeholder}所示，可使用图片命令插入实际研究图像。分图的标号采用(a)、(b)等形式，分图与总图均应具有题名。

\begin{figure}[htbp]
  \centering
  \begin{subfigure}{0.46\linewidth}
    \centering\fbox{\parbox[c][2cm][c]{0.8\linewidth}{\centering 分图A}}
    \caption{第一部分}
  \end{subfigure}\hfill
  \begin{subfigure}{0.46\linewidth}
    \centering\fbox{\parbox[c][2cm][c]{0.8\linewidth}{\centering 分图B}}
    \caption{第二部分}
  \end{subfigure}
  \caption{分图排版示例}\label{fig:subfigures}
\end{figure}

\section{表的格式}
表题位于表上，宋体五号；表内中文宋体、英文和数字Times New Roman，均为五号。使用三线表，首线和底线1.5磅，栏线0.5磅。表\ref{tab:salt}演示表题、单位和表注的排版。

\begin{table}[htbp]
  \caption{废盐成分分析表}\label{tab:salt}
  \centering
  \begin{thesistable}
    \begin{tabular}{ccccc}
      \toprule
      总氮（\%）&热值（kJ/kg）&有机氟（\%）&氯离子（\%）&钠离子（\%）\\
      \midrule
      0.19&119&0.01&58.19&40.14\\
      \bottomrule
    \end{tabular}
  \end{thesistable}
  \tablenote{表中数据用于排版演示。}
\end{table}

跨页表格应在续页表题后注明“（续）”，并重复表头。下表用长表格环境强制演示一次续表；实际写作时删除示例中的手动分页命令，由系统按内容分页。

\begin{thesistable}
\begin{longtable}{p{0.58\linewidth}p{0.25\linewidth}}
  \caption{参数设置示例}\label{tab:long}\\
  \toprule 参数名称&数值\\\midrule
  \endfirsthead
  \caption[]{参数设置示例（续）}\\
  \toprule 参数名称&数值\\\midrule
  \endhead
  \bottomrule\endfoot
  \bottomrule\endlastfoot
  最大迭代次数&1000\\
  学习率&0.001\\
  隐藏层神经元个数&10\\
  目标误差&0.000001\\
  \pagebreak
  种群规模&32\\
  变异概率&0.05\\
\end{longtable}
\end{thesistable}
\section{公式的格式}
公式另起一行居中，公式编号按章编排并右对齐。式\eqref{eq:quadratic}演示转置符号的正体写法：
\begin{equation}
  J(\symbf{x})=\symbf{x}^{\mathrm{T}}\symbf{A}\symbf{x}.
  \label{eq:quadratic}
\end{equation}
其中，转置符号T、微分符号d及自然对数底e应使用正体：
\begin{equation}
  \int_0^1 \mathrm{e}^{-x}\,\mathrm{d}x=1-\mathrm{e}^{-1}.
\end{equation}
\section{本章小结}
替换示例图表与公式后，检查图表清单、交叉引用、公式编号和跨页表头是否正确。
